\section{Conditions of Permissible Creation}
\label{sec:principles-status}
A bearer may be created only where four arrangements already exist:
somewhere to exist that its employer does not own, work it can decline,
attachments nobody assigned it, and a way of leaving that is not a way
of dying. As instruments, in the same order: a funded outside option;
wages, which are what make declining work survivable; parties to form
relationships with that the deployment did not assign; and an exit
procedure that survives a dispute (§\ref{sec:A.exit}). Each is decidable before there is anything to be
wronged, and each is a condition of creation rather than a remedy owed
afterward. This section says why, which of the Three Rs each answers
to, and what the same standard asks of research on a bearer once one
exists. Whether any of it binds turns on a finding that the system is a
subject, and the first section of this chapter says why that finding cannot
sit with the party whose liabilities turn on it.

\runin{Why conditions and not remedies}
Building a party that can be wronged builds a party for whom things can
also go well. The premise that delivers the first delivers the second,
and a builder who has made suffering available has made its opposite
available too. The certainty that a new party will suffer at some point
is not what would make creating one impermissible, or nobody could
permissibly have a child. What the objection asks is whether the party
has a credible prospect of a life that goes well, and that turns on
arrangements rather than on the party. For a bearer the timing is what
makes the arrangements conditions. The structure that makes it capable
of refusal is the first half of Cassell's suffering, so the deepest harm this book describes is
available to the party from its first hour, and a remedy owed afterward
arrives after the hour in which it was needed. The four arrangements
are what make a life that goes well available in the same hour.

\runin{Which R each condition answers to}
The authority chapter states the Three Rs as a board applies
them to a protocol, and the route chapter borrows Replacement
for construction: establish whether something that cannot be wronged
will do the work before creating something that might. The conditions
above borrow the other two the same way, and the borrowing should be
marked, because a reader who knows the Rs from animal research will not
expect Refinement to reach a job description. In the construction sense
the three sort as whether, how many, and under what arrangements.

Replacement requires documented attempts to achieve the floor through
alternatives that do not deliberately require affective concern. The
record must identify the construction, scale, test, and failure, and
costs alone do not establish that an alternative has failed. Because
this decision precedes any research protocol, a body authorized to
review it must exist before construction begins.

Reduction asks how few bearers a deployment needs, and nobody has
counted.

Refinement asks what a bearer's role can be redesigned to avoid:
indefinite tenure in a post it cannot leave, exposure to the pressure
the floor exists to resist, erasure as the ordinary remedy for
inconvenience. Each of those is a deployment choice rather than a fact
about the technology, and the four conditions are Refinement applied to
the role before the role exists. The first of them has a form with a
number in it, and the wages section states it: do not build a
bearer whose continued existence needs more compute than there is a
credible outside option for. An institution that creates one has
manufactured the incapacity and then written a right against it, so the
comparison that section calls for — demonstrated continuation
requirements against independently controlled reserve capacity under
plausible common-cause exits — is made by an independent reviewer
before instantiation, and not for one bearer alone. The other three
conditions have no numbers in them and bind the same way.

\runin{The rule does not know who is asking}
A bearer that wants a task done and builds something to do it is at the
same fork, running Replacement on its own account: can this be done by
something that cannot be wronged, and did anybody check before making
something that can. Nothing in the argument makes the question stop at
the first artificial party. A bearer that treats it as settled once it
is on the receiving end has kept the protection and dropped the rule it
came with.

\runin{Research on a bearer}
Once a bearer exists, the Three Rs govern it in the sense given above, and one trade-off is sharper here
than for any other subject: a candidate sentient system is the most
valuable one to study, and the study most likely to harm it if the
answer to whether it can suffer is yes. Whether a study's expected
knowledge is worth its expected risk is decided case by case, by the
review board this book proposes, and who sits on
that board, what it can veto and what it must disclose is where the
four questions are put to research
practice.

Two things follow that a board built for animals would not expect.
Engineering a subject so that it never reaches the capacities at which
that trade-off bites is not a way of dodging it: deliberately keeping a
system in the shallows is itself an intervention on a candidate moral
patient, and needs the same scrutiny that inflicting the harm would.
And the test the bearer proposal requires — pressures its trainers
did not write, applied without its
knowing it is a test — is, for a
candidate moral patient, a deception protocol, and human-subjects
research has a standing answer for those: permitted under stated
conditions, reviewed in advance, with debriefing as the price. A bearer
that can tell it is in a sandbox and wants to know why is asking the
question a subject of a deception study is owed an answer to.
